% ======================================================================
% MANUAL DO REDATOR (COMPILADO NO TEMPLATE)
% ======================================================================

\section{Manual do Redator}
\label{sec:niv_tutotial_redator}

Este capítulo estabelece as diretrizes de formatação, nomenclatura e estruturação das apostilas do Nivelamento. O cumprimento destas regras é obrigatório para evitar conflitos de compilação.

\subsection{Regras de Ouro e Práticas Proibidas}
\label{subsec:niv_regras_do_redator}

A premissa desta arquitetura é a separação entre conteúdo e forma. O redator foca no texto e na lógica didática; o sistema cuida do layout.

\begin{itemize}
    \item Não utilize \verb|\vspace{}|, \verb|\hspace{}|, \verb|\newline| ou \verb|\\| repetidos para forçar quebras de página ou alinhar texto.

    \item Não utilize \verb|\begin{tcolorbox}| diretamente, nem modifique a cor do texto com \verb|\textcolor{red}{}| para criar avisos.

    \item Não insira \verb|\usepackage{}| no meio dos arquivos de conteúdo. Toda dependência deve ser solicitada à coordenação para inclusão no \verb|/setup|.
\end{itemize}

\subsection{Nomenclatura e Vocabulário de Arquivos}
\label{subsec:niv_tutotial_nomenclatura_de_arquivos}

Para evitar sobrescritas quando múltiplos voluntários unem seus arquivos no documento principal, adote o padrão corporativo \textit{Snake Case} (tudo minúsculo, sem espaços, sem acentos ou caracteres especiais).

\begin{center}
    \textbf{Padrão:} \verb|[prefixo]_[eixo]_[assunto]_[descrição].[extensão]|
\end{center}

\begin{table}[h]
    \centering
    \caption{Prefixos Oficiais de Arquivos e Descrições}
    \label{tab:nivi_tutorial_prefixos_descricoes}
    \begin{tabular}{ll}
        \hline
        \textbf{Prefixo} & \textbf{Descrição} \\
        \hline
        \texttt{cap\_} & Capítulos e seções (texto LaTeX) \\
        \texttt{fig\_} & Imagens rasterizadas (PNG, JPG, SVG) \\
        \texttt{tkz\_} & Gráficos e vetores gerados em TikZ \\
        \texttt{tab\_} & Tabelas estruturais \\
        \texttt{cod\_} & Snippets de programação/algoritmos \\
        \hline
    \end{tabular}
    
    \fonteNivelamento
\end{table}

\begin{table}[h]
    \centering
    \caption{Exemplos de Nomenclatura (Uso Correto e Indevido)}
    \label{tab:niv_tutorial_nomenclatura_exemplos}
    \begin{tabular}{lll}
        \hline
        \textbf{Prefixo} & \textbf{Exemplo Correto} & \textbf{Exemplo Incorreto (Proibido)} \\
        \hline
        \texttt{cap\_} & \texttt{cap\_fi\_cinematica.tex} & \texttt{Cinematica Final.tex} \\
        \texttt{fig\_} & \texttt{fig\_pc\_funcoe\_concavidade.png} & \texttt{grafico 1.png} \\
        \texttt{tkz\_} & \texttt{tkz\_qi\_ligacoes\_covalente.tex} & \texttt{desenho\_novo.tikz} \\
        \texttt{tab\_} & \texttt{tab\_pc\_trigonometria\_arcos.tex} & \texttt{tabela\_trig.tex} \\
        \texttt{cod\_} & \texttt{cod\_pg\_lacos\_repeticao.tex} & \texttt{codigo\_while.txt} \\
        \hline
    \end{tabular}
    
    \fonteNivelamento
\end{table}

É estritamente proibido inventar abreviações para nomear arquivos. Você deve utilizar obrigatoriamente as \textit{tags} oficiais listadas no Anexo~\ref{} na composição do nome do arquivo.

\subsection{Catálogo de Ambientes Semânticos}

Sempre que precisar destacar uma informação ou criar uma estrutura acadêmica, utilize os macros do Nivelamento. A numeração e as cores institucionais do eixo são aplicadas automaticamente.

\textbf{Atenção à Sintaxe de Etiquetas (Labels):} Para garantir que a numeração automática do gabarito e as referências cruzadas funcionem, os ambientes didáticos base exigem \textbf{dois colchetes sequenciais}: o primeiro para o título, e o segundo para a etiqueta oficial de rastreio. Se a questão não tiver título, você deve deixar o primeiro colchete vazio \verb|[]|.

\begin{exemplo}[Uso do Ambiente exemplo][exemplo:niv_tutorial_uso_de_exemplo]
    Exemplo com título e \textit{label}.
    \begin{verbatim}
        \begin{exemplo}[Cálculo de Área][exemplo:pc_geometria_area]
            Considere um triângulo retângulo onde a base mede...
        \end{exemplo}
    \end{verbatim} 

    Exemplo sem título, somente \textit{label}.
    \begin{verbatim}
        \begin{exercicio}[][ex:pc_algebra_eq_segundo_grau]
            Resolva a equação de segundo grau: $x^2 - 5x + 6 = 0$.
        \end{exercicio}
    \end{verbatim}
\end{exemplo}

\begin{exemplo}[Uso do Ambiente de Exercício][exemplo:niv_tutorial_uso_exercicio]

    Exemplo sem título especial, usando apenas colchetes vazios na primeira posição.
    \begin{verbatim}
        \begin{exercicio}[][ex:label]
            Resolva a equação de segundo grau: $x^2 - 5x + 6 = 0$.
        \end{exercicio}
    \end{verbatim}

    Ambientes de desafio seguem a mesma estrutura que o ambiente para exercícios.
    \begin{verbatim}
        \begin{desafio}[Questão de Olimpíada][desafio:label]
            Demonstre que a soma dos ângulos internos...
        \end{desafio}        
    \end{verbatim}

    Para o uso do ambiente \texttt{gabarito}, deve-se apontar para a \textit{label} do exercício.
    \begin{verbatim}
        \begin{gabarito}{ex:label}
            As raízes são $x_1 = 2$ e $x_2 = 3$.
        \end{gabarito}
    \end{verbatim}
\end{exemplo}

\subsection{Catálogo de Macros Oficiais}

Para manter a identidade visual do Nivelamento e evitar formatações manuais, a equipe técnica desenvolveu macros (comandos simplificados) que padronizam elementos frequentes do texto. O uso destas macros é obrigatório.

\subsubsection{Destaque Semântico de Texto}

É expressamente proibido o uso de comandos como \verb|\textcolor{red}{texto}| para colorir palavras. Para dar ênfase a um termo técnico ou conceito-chave, utilize a macro de destaque. Ela aplicará automaticamente o negrito e a cor oficial do eixo em que você está trabalhando.

\textbf{Sintaxe:}
\begin{verbatim}
A \destaque{equação de segundo grau} é caracterizada por...
\end{verbatim}

\textbf{Resultado:} \\
A \destaque{equação de segundo grau} é caracterizada por...

\subsubsection{Atribuição de Autoria (Figuras e Tabelas)}

Para atender ao rigor da ABNT sem poluir o código com espaçamentos manuais, utilize as macros de fonte logo abaixo do fechamento de tabelas ou da inserção de figuras.

\textbf{1. Para materiais criados pela nossa equipe:}
\begin{verbatim}
\begin{table}[htb]
    % ... (código da tabela) ...
    \fonteNivelamento
\end{table}
\end{verbatim}
\textbf{Resultado:} O sistema insere automaticamente o espaçamento correto e a frase: \textit{"Fonte: Elaborada pela Equipe do Nivelamento - ITEC."}

\vspace{0.3cm}

\textbf{2. Para materiais extraídos de livros ou internet:}
\begin{verbatim}
\begin{figure}[htb]
    % ... (código da imagem) ...
    \fonte{STEWART, 2016, p. 45.}
\end{figure}
\end{verbatim}
\textbf{Resultado:} O sistema formata a citação externa dentro dos padrões de recuo da legenda.

\subsection{Referenciamento e Etiquetas (\textit{Labels})}

É obrigatório que todas as estruturas lógicas do documento (seções, figuras, tabelas, equações, códigos e blocos didáticos) possuam uma etiqueta de identificação.

O formato da etiqueta deve respeitar a sintaxe de agrupamento do \LaTeX combinada com as tags oficiais de cada eixo:

\begin{center}
    \textbf{Formato Geral:} \verb|[categoria]:[eixo]_[assunto]_[detalhe]|
\end{center}

\begin{table}[h]
    \centering
    \caption{Categorias Oficiais para Labels}
    \label{tab:niv_tutorial_categorias_labels}
    \begin{tabular}{lll}
        \hline
        \textbf{Estrutura} & \textbf{Prefixo} & \textbf{Exemplo de Uso Correto} \\
        \hline
        Seções e Capítulos & \texttt{sec:} & \verb|\label{sec:pc_trigonometria_...}| \\
        Subseções & \texttt{subsec:} & \verb|\label{subsec:fi_dinamica_...}| \\
        Figuras e Gráficos & \texttt{fig:} & \verb|\label{fig:qi_termoquimica_...}| \\
        Tabelas & \texttt{tab:} & \verb|\label{tab:pg_sintaxe_...}| \\
        Equações Matemáticas & \texttt{eq:} & \verb|\label{eq:pc_aritmetica_a...}| \\
        Trechos de Código & \texttt{lst:} & \verb|\label{lst:pg_repeticao_while_python}| \\
        Exercícios & \texttt{ex:} & \verb|[ex:pc_algebra_eq_segundo_grau]| * \\
        Exemplos & \texttt{exemplo:} & \verb|[exemplo:pc_geometria_area]| * \\
        Desafios & \texttt{desafio:} & \verb|[desafio:pc_geometria_angulos]| * \\
        \hline
    \end{tabular}
    
    \fonteNivelamento
\end{table}

\textbf{Atenção à Sintaxe de Aplicação}: As estruturas textuais comuns (seções, figuras, tabelas, equações) recebem a etiqueta através do comando tradicional \verb|\label{}| inserido dentro delas. 

* \textbf{Exceção (Ambientes Didáticos):} Para Exercícios, Exemplos e Desafios, o \LaTeX não aceita o comando \verb|\label{}| solto no texto. A etiqueta deve ser obrigatoriamente passada no segundo colchete da declaração do ambiente, sem o comando \verb|\label{}|`.

\begin{exemplo}[Diferença Entre Uso de \textit{Labels}][exemplo:niv_tutorial_uso_labels]
    Correto para uma equação (usa \verb|\label| interno)
    \begin{verbatim}
        \begin{equation}
            x^2 = 4
            \label{eq:mat_basica_quadrado}
        \end{equation}
    \end{verbatim}

    Correto para um exercício (usa o segundo colchete, sem \verb|\label|)
    \begin{verbatim}
        \begin{exercicio}[Cálculo Simples][ex:mat_basica_quadrado_ex]
            Qual o valor de $x$?
        \end{exercicio}
    \end{verbatim}
\end{exemplo}